\section{CAF Origins and Activation Mechanisms}
\label{sec:origins}

The cellular origins of CAFs in breast cancer are more diverse than previously appreciated. While tissue-resident fibroblasts were long considered the primary source, accumulating evidence indicates that CAFs can arise from multiple precursor populations through distinct activation pathways.

\subsection{Tissue-resident fibroblasts as the predominant source}

In the normal mammary gland, fibroblasts exist in a quiescent state, maintaining tissue homeostasis through controlled ECM production and paracrine signalling with epithelial cells \citep{plaks2015mammary}. Upon malignant transformation, these resident fibroblasts become activated through exposure to tumour-derived factors, most notably transforming growth factor-$\beta$ (TGF-$\beta$), platelet-derived growth factor (PDGF), and fibroblast growth factors (FGFs) \citep{kalluri2016biology}.

TGF-$\beta$ signalling represents the best-characterised activation pathway. Tumour cells and immune cells secrete latent TGF-$\beta$, which is activated in the ECM through integrin-mediated mechanical force or proteolytic cleavage \citep{shi2011latent}. Active TGF-$\beta$ binds to TGF-$\beta$ receptors on fibroblasts, triggering SMAD2/3 phosphorylation and nuclear translocation, which drives the transcription of myofibroblast markers including $\alpha$-smooth muscle actin ($\alpha$-SMA), collagen type I, and fibronectin \citep{derynck2001smad}. This TGF-$\beta$-driven activation programme is particularly prominent in myCAFs, the most abundant CAF subpopulation in breast tumours \citep{erve2024}.

PDGF signalling operates in parallel, with tumour-secreted PDGF-AA and PDGF-BB activating PDGF receptor $\alpha$ (PDGFR$\alpha$) and PDGFR$\beta$ on fibroblasts, respectively \citep{ostman2001pdgf}. PDGF signalling promotes fibroblast proliferation, migration, and ECM production, and is essential for establishing the desmoplastic reaction characteristic of many breast cancers \citep{shao2016osteonectin}.

\subsection{Bone marrow-derived mesenchymal stem cells}

A proportion of CAFs originate from bone marrow-derived mesenchymal stem cells (BM-MSCs) that are recruited to the tumour site through chemokine gradients \citep{spaeth2009inflammation}. CXCL12 (SDF-1) secreted by tumour cells and activated stroma attracts CXCR4-expressing BM-MSCs, which subsequently differentiate into CAFs under the influence of tumour-derived TGF-$\beta$ and IL-6 \citep{karnoub2007mesenchymal}.

In breast cancer, BM-MSC-derived CAFs have been shown to promote tumour growth, angiogenesis, and metastasis more potently than their tissue-resident counterparts \citep{quante2011bone}. This functional superiority may reflect the acquisition of epigenetic modifications during bone marrow residence that prime these cells for rapid activation upon tumour recruitment. However, the relative contribution of BM-MSCs to the total CAF pool varies across breast cancer subtypes and individual tumours, with estimates ranging from 5\% to 25\% \citep{augello2005bone}.

\subsection{Endothelial-to-mesenchymal transition}

Endothelial cells can acquire fibroblast-like properties through endothelial-to-mesenchymal transition (EndMT), a process driven by TGF-$\beta$, Wnt, and Notch signalling \citep{zeisberg2007endothelial}. During EndMT, endothelial cells downregulate vascular endothelial cadherin (VE-cadherin) and CD31 while upregulating $\alpha$-SMA, fibroblast-specific protein 1 (FSP1), and collagen production \citep{potenta2008origins}.

In breast cancer, EndMT-derived CAFs have been identified in the perivascular niche, where they contribute to vascular normalization and immune cell trafficking \citep{gasparri1999thrombospondin}. These cells may represent a distinct CAF subpopulation with specialized functions in angiogenesis regulation and immune cell extravasation, though their precise contribution to the overall CAF landscape remains to be fully characterized.

\subsection{Epithelial-to-mesenchymal transition as a source}

Epithelial-to-mesenchymal transition (EMT) can generate CAF-like cells from malignant epithelial cells themselves \citep{petersen2003epithelial}. Through EMT, tumour cells lose epithelial markers (E-cadherin, cytokeratins) and acquire mesenchymal properties (N-cadherin, vimentin, $\alpha$-SMA), becoming indistinguishable from fibroblasts by standard markers \citep{kalluri2009emt}.

This phenomenon---termed ``tumour-cell-derived CAFs''---is particularly relevant in breast cancer, where EMT is a central mechanism of metastasis and therapeutic resistance \citep{niegel2018}. However, distinguishing tumour-derived CAFs from bona fide fibroblast-derived CAFs requires lineage tracing or tumour-specific genetic markers, which are rarely available in clinical samples. The functional significance of tumour-derived CAFs remains debated: some studies suggest they are more aggressive than conventional CAFs \citep{chen2014}, while others argue they represent a minor population with limited impact on tumour biology \citep{strell2019}.

\subsection{Activation signalling networks}

Beyond the canonical pathways described above, CAF activation in breast cancer involves a complex network of interacting signals:

\textbf{IL-6/STAT3 axis:} Tumour-derived IL-6 activates STAT3 signalling in fibroblasts, promoting their differentiation into iCAFs characterized by high expression of inflammatory cytokines (IL-6, IL-8, CXCL1) and chemokines (CCL2, CCL5) \citep{erve2024}. This axis creates a feed-forward loop: iCAF-secreted IL-6 further activates STAT3 in tumour cells, enhancing their proliferative and invasive capacity.

\textbf{Hedgehog signalling:} Tumour-secreted Sonic hedgehog (SHH) activates the Hedgehog pathway in fibroblasts through the receptor PTCH1 and effector GLI2, driving ECM production and desmoplasia \citep{rao2006hedgehog}. In breast cancer, Hedgehog signalling is particularly active in the basal-like subtype, where it contributes to the characteristic dense stroma.

\textbf{Notch signalling:} Direct cell--cell contact between tumour cells and fibroblasts activates Notch signalling in the latter, promoting their proliferation and differentiation \citep{strell2019}. This juxtacrine mechanism operates independently of secreted factors and may be particularly important at the tumour--stroma interface.

\textbf{Mechanical signalling:} The stiff, collagen-rich ECM produced by activated fibroblasts generates mechanical forces that are sensed by mechanosensitive ion channels and integrins on fibroblasts themselves, reinforcing their activated phenotype \citep{calvo2013}. This mechanical feedback loop contributes to the self-sustaining nature of the CAF phenotype once established.

\subsection{Heterogeneity of activation states}

Importantly, fibroblast activation is not a binary switch but a continuum of states. Single-cell studies have revealed that CAFs within a single tumour occupy multiple activation states, ranging from quiescent to fully activated \citep{erve2024, bartoschek2018spatially}. The balance between these states is influenced by local signalling gradients, ECM composition, and oxygen tension.

This heterogeneity of activation states has important implications for therapeutic targeting: interventions that block a single activation pathway may shift the equilibrium toward alternative states rather than eliminating CAFs entirely. A systems-level understanding of the activation network is therefore essential for designing effective CAF-targeted strategies.
